\section{Conclusions}
\label{sec:conclusions}

\begin{itemize}
    \item P3Net combines P3's dependency modelling (a linkage tree built from pairwise mutual information over the genotype) with a relative surrogate that screens candidates before real evaluation, testing whether this combination improves surrogate-loop efficiency over NSGANetV2~\citep{lu2020nsganetv2} on a shared cell search space extended to joint architecture-and-hyperparameter NAS. The answer is no: P3Net does not separate from \texttt{random\_search} in any of the sixteen (search-space, budget) cells measured (Results), and its margin against every differently-engineered baseline correlates almost perfectly with that baseline's own margin against \texttt{random\_search} ($\rho=0.962$, $n=80$) -- indicating this comparison is dominated by a mechanism shared across arms, not a property specific to the relative surrogate. (The wider, ten-baseline version of this same correlation, run after the ablation family below was added to the grid, is reported further down.) The mechanistic account for this result is developed below.
    \item Limitations: surrogate quality plausibly depends on genotype encoding quality, though this dependency is not isolated by a dedicated ablation in the present design; the linkage tree is built from pairwise mutual information, which may be too coarse a signal to capture higher-order, deeply structured dependencies. JAHS-Bench-201 answers every query through a learned surrogate rather than measured data, so results on it inherit that surrogate's approximation error relative to real, independent training; since every compared method queries the same surrogate, systematic bias affects them equally, but this does not rule out the error correlating specifically with the structure P3Net exploits. A result that also holds on NAS-HPO-Bench-II, built from measured data, is therefore treated as the primary finding, and a result specific to JAHS-Bench-201 alone as requiring confirmation. NAS-HPO-Bench-II is inherently CIFAR-10-only; JAHS-Bench-201's single-dataset limitation raised by an earlier pass of this project (and shared with the closest prior work~\citep{bartnik2026evolutionary}, Related Work) is addressed by extending the grid to run across all three of JAHS-Bench-201's own built-in datasets (CIFAR-10, Colorectal-Histology, Fashion-MNIST) rather than only its CIFAR-10 default. This checks robustness across that benchmark's own dataset family rather than generalisation beyond it, and the honest result of that check is heterogeneity, not confirmation: the core surrogate-vs.-no-surrogate mechanism holds most completely on Fashion-MNIST, partially on CIFAR-10, and narrows on Colorectal-Histology (Results, ``Scope, and what would extend it''). NAS-HPO-Bench-II's own single-dataset ceiling still bounds what the measured-data leg of the comparison can speak to, and on that benchmark specifically the ablation result disappears and inverts at the highest budget tested. Results are further scoped to a single, shared architecture family (a NAS-Bench-201-style cell search space) and small-scale image classification across both benchmarks; neither a different architecture family nor a different task domain is tested here, and either would need different benchmark infrastructure entirely, not a configuration change. This specific limitation is not unique to the present comparison: the wider NAS benchmarking literature documents that both search-space design and task domain materially affect method rankings, and that conclusions drawn on one NAS benchmark frequently fail to transfer to another~\citep{yang2020nasfrustratinglyhard,mehta2022nasbenchsuite,lopes2023nasbenchdesign,tu2022nasbench360,duan2021transnasbench101} (Related Work).
    \item Future directions: stronger surrogate model (GNN over the cell graph), active learning to decide which candidates to evaluate fully, extension to architecture spaces with variable depth; training the surrogate on lower-fidelity observations $f_1^{(k)}$, $k<K$, rather than exclusively on full evaluations at $r_K$ as in the present work, using the fidelity ladder already introduced in Section~\ref{sec:problem-formulation}. This is left out here to keep the search loop's cost reduction cleanly attributable to the relative surrogate alone, but is a natural next lever once that attribution is established. Testing beyond this benchmark family's shared architecture family or task domain, or replacing benchmark-substrate evaluation with live training, would need substantially different infrastructure and is left as a separate undertaking. This is not merely a generic call for more testing: two specific observations from the present results motivate it directly (Results, ``Scope, and what would extend it''). The relative surrogate's measured advantage over the matched ablations, and the order-of-magnitude gap it opens in genotype duplication rate, both hold on a search space the wider NAS literature documents as comparatively small and unstructured~\citep{lopes2023nasbenchdesign}, while the P3/GOMEA lineage this work builds on has repeatedly shown that its advantage over generic search grows, rather than shrinks, as problem size and structural complexity increase~\citep{goldman2015fast,qiao2024gomearuntime,dushatskiy2019csgomea} (Related Work). That combination is a reasoned, evidence-based motivation for testing the same mechanism at larger scale or in a different task domain -- not a demonstration that it will succeed there, which only running that comparison could show. Extending to a new task domain is itself a precedented move in this literature, not a leap into the unknown: MoSegNAS carries the NSGANetV2~\citep{lu2020nsganetv2} lineage into semantic segmentation, but only after redesigning both the surrogate and the search space for the new task~\citep{lu2022mosegnas}, which is exactly why the present work treats a domain change, like live training, as a separate undertaking with its own infrastructure requirements rather than a configuration change on top of the present grid.
    \item A design-decision ablation round (Results, ``Design decisions: a P3Net ablation round'') resolves several open points empirically. Donor provenance and the population-truncation rule show no measurable effect across the full grid and are kept as designed. Computing $f_2$ fresh via the substrate's analytic cost hook, rather than inheriting it from the ancestor, gives a consistent, partly-significant hypervolume improvement and is adopted as the default -- though it also lowers acceptance-gate precision (Results, Surrogate quality); a follow-up investigation rules out a measurement artefact, accumulated chain error, a worse-fitting surrogate, and extrapolation distance as the cause without fully resolving why a lower-precision gate coexists with a better downstream result, left as a target for follow-up surrogate diagnostics. A simultaneous-cascade variant of the pyramid (every level swept every call, not only the newest) gives a similarly mixed, dataset-dependent result and a real per-iteration compute cost, and is not adopted.
    \item A later investigation into the per-dataset ranking pattern (Results, ``Diagnostics: population-pyramid bootstrap share'') identifies a more consequential mechanism: \texttt{Pyramid}'s geometric level growth, combined with starting every new level from an empty, uniform-randomly-bootstrapped population and sweeping only the newest level, structurally forces the uniform-random share of P3Net's own evaluation budget to rise with budget (measured $81.6\%\to93.0\%$ from budget~50 to~350 on NAS-HPO-Bench-II, cross-validated at $93$--$94\%$ across all four search spaces at budget~350). This share is large enough to explain why P3Net is not statistically distinguishable from \texttt{random\_search} in any of the sixteen cells measured, and why its margin against every differently-engineered baseline tracks that baseline's own margin against \texttt{random\_search} almost exactly ($\rho=0.962$, $n=80$): a substantial share of both the budget trend (widening losses to SH-EMOA/TPE/NSGA-Net/NSGANetV2) and the dataset trend (smallest losses on Fashion-MNIST, the search space where guided search barely beats uniform sampling at all) documented above is attributable to this pyramid-growth artefact rather than to a comparative property of the relative surrogate itself. A joint $\kappa$/acceptance-threshold sensitivity sweep (Results, ``Surrogate error accumulation'') rules out retuning either knob as a fix: the bootstrap-dominated share and fixed-budget hypervolume are both essentially flat across the full sensitivity grid.

    A redesign that warm-starts each newly grown level from $\mathcal{H}_t$'s current nondominated front,
rather than an empty population, and replaces the single-pass strict-Pareto-dominance stall trigger with a
hypervolume-contribution criterion, reduces the bootstrap-dominated share to $90$--$92\%$ across all four
search spaces at budget~350 (down from the $93$--$94\%$ documented above), confirmed on the full $R=30$ grid.
This is a real, directionally correct improvement, though a more modest one than the severity of the original
finding might suggest: the number of distinct levels grown by budget~350 is unchanged (7, sizes
$\{2,4,8,16,32,64,128\}$), meaning each level survives somewhat longer without meaningfully outpacing the
geometric progression's own growth rate. The consequence for the headline comparison: P3Net still does not
separate significantly from \texttt{random\_search} in any of the sixteen cells even under the fixed pyramid
($0/16$, Holm-corrected) -- the mechanism identified is real and the fix measurably shrinks it, but not by
enough, on this benchmark family and scale, to flip the headline finding. A further attempt -- starting each
level's real mixing much sooner, closer to how canonical P3~\citep{goldman2014parameterless} generates and
cascades one individual at a time rather than bootstrapping an entire level's target size with pure random
draws first -- causes a severe regression instead: this project's donor-pool-dependent optimal mixing, unlike
canonical P3's~\citep{goldman2014parameterless} single-individual hill climbing, degenerates without a real population to mix from. A
donor-pool-aware version of that same idea, not a hand-picked constant tuned against these four benchmarks
specifically, is left as the natural next candidate.
    \item A related but structurally independent limitation, surfaced alongside the pyramid finding above, concerns $\hat{\delta}_F$'s own model class rather than the search loop around it (Results, ``Diagnostics: surrogate representational capacity''): it, and the absolute regressor shared by the P3+absolute-regressor ablation and NSGANetV2~\citep{lu2020nsganetv2}, are all a plain \texttt{sklearn.LinearRegression} over one-hot genotype features, which can only represent additive, no-interaction fitness differences -- in tension with linkage-guided crossover's own premise that the variable groups it protects have a joint, not merely additive, effect on fitness (Related Work). A measured, dataset-ordered proxy for how much this misspecification actually costs (``architecture-involving epistasis'') tracks the same per-dataset pattern the pyramid finding above already partly explains, offered as a second, independently-motivated account rather than a competing one; which of the two, or what combination, is doing the explaining was not tested interventionally. Unlike P3Net's own linear surrogate, the two closest related methods sidestep this specific failure mode by construction -- Bartnik's~\citep{bartnik2026evolutionary} absolute regressor is fit as one of several nonlinear-capable model families, and eLyMPuS~\citep{przewozniczek2026lympus} is not a regression over a flat encoding at all -- making this a concrete, literature-motivated candidate for a follow-up alongside the pyramid redesign above, not a hypothetical one. A candidate fix, scoped to keep the relative/delta formulation's own data-efficiency advantage rather than swapping to a heavier model class, augments $\hat{\delta}_F$ with binary ``did coordinates $i$ and $j$ change together'' interaction columns, one per unordered coordinate pair, self-scoping to whichever linkage subset a given pair actually differs within without any separate bookkeeping. On a constructed non-additive pattern, unrepresentable by the original additive-only encoding by construction, the augmented surrogate fits it near-exactly -- direct evidence the added capacity is real. On the full $R=30$ grid as the new arm \texttt{p3net\_surrogate\_interactions}, the result is an honest null: against the Phase-2-fixed \texttt{p3net} baseline, significance is reached in $0/16$ cells, with small and sign-inconsistent effect sizes. The representational limit is real and measurably fixed, but this benchmark family's landscapes evidently do not carry enough exploitable interaction structure, at this scale, for the added capacity to change search outcomes -- consistent with the "not enough architectural complexity yet" account of P3Net's weak overall showing discussed above (Related Work).
    \item A third, complementary finding (Results, ``Surrogate quality: magnitude calibration'') concerns numerical stability rather than representational capacity or search-loop mechanics: $\hat{\delta}_F$'s predicted-delta \emph{magnitude} is calibrated far worse than trivially predicting the mean true delta in every search space, concentrated overwhelmingly in one specific pyramid level's first sweep pass immediately after its own bootstrap completes (a $6$--$19\times$ higher median squared error than any other level size, on real, repeated measurement). This is the textbook failure mode a point-estimate surrogate trusted without regard to its own reliability exhibits under a small, evolving training sample, for which the surrogate-assisted evolutionary computation literature's established remedy is regularisation or explicit uncertainty-aware model management~\citep{jin2011surrogate,hoerl1970ridge} -- distinct from, and adopted alongside, the pyramid-transition redesign above, since both concern exactly that same moment. Adopting \texttt{RidgeCV} as $\hat\delta_F$'s regressor does not correct it: on the full $R=30$ grid at budget~350, level~16 remains by far the worst-calibrated level ($R^2=-4.71$, MSE $766.7$) against every other level (best: level~2, $R^2=+0.55$, MSE $85.2$), a $\approx 9\times$ gap squarely inside the originally-documented $6$--$19\times$ range. Regularisation is the textbook remedy for unstable small-sample linear regression, but it does not, on this evidence, correct a systematic miscalibration concentrated at one specific, structurally-determined pyramid level -- the two problems this literature conflates (unstable coefficients vs.\ a landscape genuinely harder to model at that particular $|\mathcal{H}_t|$) may not be the same problem here, left as an open point for follow-up surrogate diagnostics rather than resolved by this pass.
    \item The P3+absolute-regressor ablation's own confound, noted above and in Results (``Diagnostics: surrogate representational capacity''), is closed by sharing P3Net's own \texttt{Pyramid} engine -- including the warm-start and hypervolume-contribution changes described above -- rather than a separate flat, fixed-size population, so this ablation isolates the surrogate alone, as the Baselines paragraph claims. On the full $R=30$ grid under the shared engine, the result is a null: P3Net vs.\ P3+absolute-regressor reaches significance in $0/16$ cells, in either direction. Once the population-engine confound that made the original comparison unfair is removed, this grid shows no measurable advantage of the relative surrogate over the absolute regressor -- a materially different, more cautious conclusion than the Baselines paragraph's original framing assumed, and one that should be read together with the other three null results just above: on this benchmark family and scale, none of the surrogate-side interventions tested in this pass (interaction features, RidgeCV, engine-sharing) separated from their respective baselines, even though each is independently, mechanistically confirmed to do what it was designed to do. The broader, every-arm-vs-\texttt{random\_search} check (Results, ``Diagnostics: population-pyramid bootstrap share'') updates accordingly: the overall reject rate drops from $67/160$ ($42\%$) pre-fix to $54/160$ ($34\%$) post-fix, and the cross-baseline-margin correlation stays essentially as strong ($\rho=0.945$, $n=144$, versus the pre-fix $\rho=0.962$) -- but the single largest mover in that table is P3+absolute-regressor itself, collapsing from $10/16$ significant wins over \texttt{random\_search} under its pre-fix flat population to $0/16$ once it shares P3Net's own bootstrap-dominated engine. That collapse is direct, mechanistic evidence for the account above: P3+absolute-regressor's apparent pre-fix strength was substantially an artefact of its \emph{different, non-Pyramid} population management, not of the absolute-regressor surrogate the ablation is meant to isolate -- the same pyramid-growth mechanism responsible for P3Net's own $0/16$ record, now shown to travel with the engine, not with either surrogate.
    \item \textbf{Why the results are what they are: a unified account, and an open, undertested hypothesis.} Four findings above are mechanistically independent but jointly consistent, and worth stating together rather than only as a list of separate diagnostics. (i)~\texttt{Pyramid}'s geometric growth structurally forces an ever-larger share of the budget into uniform-random bootstrap, a real and partially-fixed defect of the \emph{search engine}. (ii)~$\hat\delta_F$ is provably additive-only, a real and partially-fixed defect of the \emph{surrogate's model class}. (iii)~The benchmark family itself -- a NAS-Bench-201-style cell space -- is independently documented in the wider NAS literature as comparatively small, redundant, and unstructured, with most randomly sampled architectures already close to optimal~\citep{lopes2023nasbenchdesign}. None of (i)--(iii) alone fully explains the result: fixing (i) and (ii) each moved the needle measurably without flipping significance, exactly what a genuinely small-landscape ceiling on (iii) would predict regardless of how well (i) and (ii) are engineered around it. The P3/GOMEA lineage's own literature is consistent with this reading: its measured advantage over generic search grows, not shrinks, with problem size and structural complexity~\citep{goldman2015fast,qiao2024gomearuntime,dushatskiy2019csgomea} -- so a small-landscape ceiling is exactly the regime in which this family is expected to show the least of what it can do, independent of any implementation defect.

    A fourth candidate, distinct from (i)--(iii) and not yet disentangled from them, follows directly from the comparison to the closest prior work. Bartnik's SA-P3-GOMEA already demonstrates that a surrogate-assisted linkage learning algorithm \emph{can} win on a NAS-Bench-201-scale search space~\citep{bartnik2026evolutionary} -- an architecture-only one, with no hyperparameter dimension at all, using an absolute, nonlinear-capable regressor, and compared against MO-LS/MO-P3-GOMEA rather than NSGA-II/NSGANetV2. The present work differs from that success on all three axes simultaneously (surrogate formulation, baseline set, and, centrally, the joint architecture-and-hyperparameter genotype this work specifically contributes), so none of the results above can attribute the gap to any one of them in isolation. The joint-genotype axis is the one this work's own contribution rests on, and also the one never isolated: extending the genotype with discretised training hyperparameters ($\Theta$) grows $n$ and, with it, the chain-depth bound $\kappa=2\lceil\log_2(n)\rceil$ and the linkage tree's own search space of candidate subsets, but training hyperparameters such as learning rate or weight decay plausibly carry less multi-way, combinatorial epistasis than architectural edge choices do -- diluting, rather than enriching, the fraction of the linkage tree's sweep budget spent on subsets with real joint structure to exploit. This work's own measurements are consistent with, without confirming, that account: $98\%$ of real P3Net moves touch an architecture edge (Results, ``Diagnostics: surrogate representational capacity''), and the held-out ``architecture-involving epistasis'' proxy used there is defined specifically around architecture edges because that is where the signal concentrates -- both observations are compatible with $\Theta$ coordinates contributing comparatively little exploitable structure of the kind P3's linkage tree is built to find, though this was never tested by holding the joint-vs-architecture-only axis fixed while varying the others, and remains a hypothesis, not a finding. Section~\ref{sec:results}'s Scope discussion already treats testing at larger scale as the natural next step for (i)--(iii); isolating this fourth, joint-genotype-specific axis needs a different comparison -- P3Net run on an architecture-only search space directly comparable to Bartnik's~\citep{bartnik2026evolutionary} own, holding the surrogate and baseline set fixed.

    The isolating comparison motivated above was run as follows. \texttt{p3net} and \texttt{random\_search} were compared on architecture-only NAS-Bench-201 (no $\Theta$ at all, $n=6$ edge coordinates, $f_1=100-\text{test-accuracy}$, $f_2=\text{FLOPs}$ -- a documented deviation from Bartnik's~\citep{bartnik2026evolutionary} own measured GPU energy, which \texttt{nats\_bench}~\citep{dong2021natsbench} does not expose as a queryable field), holding the P3Net engine, surrogate, and baseline exactly as used throughout this work, at $R=30$ seeds and budgets $\{100,350\}$. The result does not support the fourth hypothesis: neither budget shows significance (adjusted $p=0.95$ at budget~$100$, $p=0.74$ at budget~$350$), and Cliff's delta is small and positive at both ($+0.02$, $+0.12$), with both methods' median relative hypervolume already at or above $0.985$. Removing the joint-genotype axis entirely -- the one candidate explanation this work's own contribution rests on -- leaves the headline result unchanged: P3Net still does not separate from random search, even on the exact class of search space where the closest prior work found separation. This shifts relative weight away from the fourth hypothesis and back toward (i)--(iii), particularly the benchmark family's own documented small-landscape ceiling~\citep{lopes2023nasbenchdesign}: both methods' near-ceiling medians on this run are consistent with a search space too easy, and too shallow in exploitable structure, for any of the search-side or surrogate-side differences among these methods to matter at this scale -- independent of whether the genotype carries a hyperparameter dimension or not.

    With the joint-genotype axis now ruled out, the remaining open explanations are the ones this account
    already narrows to: (iii) a benchmark family too small and unstructured for any tested method to separate
    from random search at this scale, and, within (i)--(ii), specific engineering choices this work's own
    \texttt{P3Net} makes differently from the two closest surrogate-assisted linkage-learning designs in the
    literature. Two concrete, direct tests follow, neither yet attempted here. First, testing (iii) properly
    means running on a larger or more structurally complex search space than any benchmark used in this
    work -- the step already motivated above (Results, ``Scope, and what would extend it''), now further
    supported by the isolation result: if the ceiling is the benchmark family itself, not the joint-genotype
    extension, then a bigger, deeper search space is the one lever left untried that this work's own diagnostics
    point to. Second, testing whether (i)--(ii) rather than (iii) are the operative bottleneck means holding this
    project's search space and evaluation infrastructure fixed while swapping in the two engines closest to
    \texttt{P3Net} in the literature it builds on: Bartnik's own \texttt{SA-P3-GOMEA}-class
    engine~\citep{bartnik2026evolutionary,dushatskiy2021novelsurrogateassistedevolutionaryalgorithm} (single-individual
    pyramid climb, nonlinear-capable absolute regressor) and Przewoźniczek et al.'s \texttt{eLyMPuS}/\texttt{OLyMPuS}
    family~\citep{przewozniczek2026lympus} (relative, non-monotonicity-based partial comparisons with guaranteed
    recursive linkage discovery, never previously evaluated on a NAS benchmark). If either reproduces this
    project's own null result on the same architecture-only search space, that is direct evidence for (iii);
    if either separates from random search where \texttt{P3Net} does not, that is direct evidence the gap sits in
    (i)--(ii), specifically in how \texttt{P3Net} instantiates them. Both comparisons are left as separate,
    dedicated undertakings rather than a configuration change on top of the present grid -- the second in
    particular requires generalising \texttt{eLyMPuS}'s binary-flip partial-comparison mechanism to NAS-Bench-201's
    five-valued categorical edges, which this work has not done and does not assume is straightforward.
\end{itemize}
